\chapter{Princípios e Fundamentos}
\label{chp:principios}

O \ppccurso da Universidade Federal de Uberlândia é orientado por princípios que articulam a excelência técnica com a responsabilidade social, a aplicação prática com a criticidade, e a formação profissionalizante com a ética. Tais princípios refletem os compromissos institucionais da UFU, as Diretrizes Curriculares Nacionais Gerais para a Educação Profissional e Tecnológica~\cite{cne_cp_1_2021} e os desafios da formação de tecnólogos(as) para a indústria automatizada do século XXI.

Entre os fundamentos que estruturam o presente Projeto Pedagógico do Curso, destacam-se:

\begin{itemize}
    \item Formação tecnológica aplicada e profissionalizante, com domínio técnico e prático das tecnologias de automação industrial, articulado à capacidade crítica de compreender e intervir no ambiente produtivo de forma ética e responsável;
    \item Integração entre fundamentos científicos e prática profissional, com abordagem sistêmica dos processos de automação — da instrumentação e do controle até a integração em redes industriais e sistemas inteligentes;
    \item Articulação entre ensino, pesquisa aplicada e extensão, como indissociabilidade fundadora da universidade pública, promovendo a formação integral do(a) estudante em contato com problemas reais e demandas do setor produtivo;
    \item Compromisso com a inovação, com incentivo à criatividade, à autonomia técnica, à cultura empreendedora e ao uso de tecnologias emergentes — como robótica e inteligência artificial aplicada — como instrumentos de modernização industrial;
    \item Inclusão e diversidade, compreendidas como dimensões fundamentais da formação acadêmica, que exigem o acolhimento de múltiplas trajetórias, saberes e identidades, bem como a construção de ambientes educativos plurais, acessíveis e antidiscriminatórios;
    \item Flexibilidade curricular e formação orientada por competências, com foco em resultados de aprendizagem diretamente aplicáveis ao exercício profissional, conforme os princípios da Educação Profissional e Tecnológica~\cite{cne_cp_1_2021};
    \item Progressividade e flexibilidade de trajetória interna, expressas na organização do curso em Área Básica de Ingresso (ABI) com o curso de Engenharia de Automação, Robótica e Inteligência Artificial (\autoref{chp:abi}): uma etapa inicial comum às duas graduações, seguida de escolha informada e autônoma do estudante pela trajetória que pretende concluir primeiro, com aproveitamento responsável dos estudos já realizados caso opte, posteriormente, por também concluir a outra formação — sempre preservando os requisitos e a identidade profissional próprios de cada graduação; distinta da mobilidade acadêmica interinstitucional tratada na Seção~\ref{sec:mobilidade};
    \item Formação orientada pela sustentabilidade, ética e responsabilidade social, com foco na eficiência energética, no uso racional de recursos e na segurança de processos industriais automatizados, capacitando o(a) estudante a integrar dimensões técnicas, sociais, ambientais e econômicas em sua atuação profissional e cidadã.
\end{itemize}

Esses princípios orientam as decisões curriculares, metodológicas e avaliativas do curso, guiando a formação de tecnólogas e tecnólogos em Automação Industrial capazes de atuar de forma crítica, aplicada e transformadora nos diversos contextos em que a automação industrial se insere.

\section{A UFU e o Curso}

A Universidade Federal de Uberlândia (UFU) estabelece, como fundamentos institucionais, a indissociabilidade entre ensino, pesquisa e extensão, o compromisso com a formação cidadã e crítica, e a valorização do conhecimento como instrumento de transformação social. Sua missão institucional é definida nos seguintes termos:

\begin{displayquote}
\textit{Desenvolver o ensino, a pesquisa e a extensão de forma integrada, realizando a função de produzir e disseminar as ciências, as tecnologias, as inovações, as culturas e as artes, e de formar cidadãos críticos e comprometidos com a ética, a democracia e a transformação social.}
\end{displayquote}

A visão que orienta sua atuação projeta a instituição como:

\begin{displayquote}
\textit{Ser referência regional, nacional e internacional de universidade pública na promoção do ensino, da pesquisa e da extensão em todos os campi, comprometida com a garantia dos direitos fundamentais e com o desenvolvimento regional.}
\end{displayquote}

O \ppccurso compartilha dessa missão e contribui ativamente para sua realização. Sua proposta pedagógica reflete os valores centrais da UFU, como o compromisso com a excelência acadêmica, a equidade, a responsabilidade social, a inovação e a formação de profissionais comprometidos com o desenvolvimento sustentável e com a transformação da realidade produtiva em que estão inseridos.

A construção deste curso busca ser conduzida de forma democrática e participativa, envolvendo docentes, técnicos-administrativos, discentes e representantes da comunidade acadêmica e regional. As decisões curriculares e metodológicas devem refletir essa escuta ampliada e estar em consonância com os planos institucionais da universidade, como o Plano Institucional de Desenvolvimento e Expansão (PIDE)~\cite{ufu_consun_31_2022} e as diretrizes do CONSUN e do CONGRAD.

Há, portanto, plena coerência entre os princípios do curso e os ideais da UFU. A formação ofertada pelo \ppccurso articula conhecimentos técnico-aplicados com o compromisso ético e social, contribuindo para a consolidação da universidade como espaço de produção de conhecimento crítico e transformador, também no campo da educação profissional e tecnológica.
